\chapter{变分原理与极值定理}

\section{泛函变分与泛函导数}

我们早已指出，算子 $(1 / i)(\partial /\partial \varphi)$ 和 $L^2$ 是厄米的，且它们的特征函数构成完备正交系。下面我们将利用更强有力的方法来证明解释这个事实的定理。

设$H$为厄米算子，并假定被$H$作用的所有函数都属于容许函数类。引进下面的算子：

\begin{equation*}
\lambda (\varphi) \equiv \frac {\int \varphi^ {*} H \varphi d \tau}{\int \varphi^ {*} \varphi d \tau}\tag{4.1}
\end{equation*}

(这里函数 $\varphi$ 可以是复的)。这个算子在每一个函数 $\varphi$ 和某个确定的数 $\lambda$ 之间建立了对应关系。由于这个理由， $\lambda(\varphi)$ 就称为泛函数。一般说来，如果指出了在函数(定义域)与数之间建立对应关系的方法，泛函数就认为是被给定了。

假定函数 $\varphi$ 变为与原来函数相邻近\textcircled{1}的新函数 $\varphi +\delta \varphi$ (当然， $\varphi + \delta\varphi$ 也必须属于容许函数类)。函数的变化可看作是在以容许类中的函数为元素的空间中的无限小转移。现在要弄清楚， $\varphi$ 作 $\delta\varphi$ 的变化，是否会引起 $\lambda$ 作相应的变化？如果是的，则要建立这些变化之间的关系。引进泛函变分 $\delta\lambda$ 的概念。设

\begin{equation*}
\begin{array}{c} \delta \int \varphi^ {*} H \varphi d \tau = \int (\varphi + \delta \varphi) ^ {*} H (\varphi + \delta \varphi) d \tau - \\ - \int \varphi^ {*} H \varphi d \tau = \int [ \varphi^ {*} H (\delta \varphi) + (\delta \varphi) ^ {*} H \varphi ] d \tau , \end{array}\tag{4.2}
\end{equation*}

这里在写出的最后表达式中，我们略去了二阶变分，例如 $\delta \varphi^{*}H\delta \varphi$ 。类似地

\begin{equation*}
\delta \int \varphi^ {*} \varphi d \tau = \int [ \varphi^ {*} \delta \varphi + (\delta \varphi^ {*}) \varphi ] d \tau .\tag{4.3}
\end{equation*}

引进记号
\begin{equation*}
A = \int \varphi^ {*} H \varphi d \tau , \quad B = \int \varphi^ {*} \varphi d \tau ,
\end{equation*}

则可写
\begin{equation*}
\lambda (\varphi) = \frac {A}{B}.\tag{4.4}
\end{equation*}

泛函数 $\lambda$ 的变分现在可以用下面的形式表示：

\begin{equation*}
\delta \lambda \equiv \frac {1}{B ^ {2}} (B \delta A - A \delta B).\tag{4.5}
\end{equation*}

(4.5)式类似于对两函数之比的导数的表达式。将它作如下的变换：

\begin{equation*}
\delta \lambda = \frac {1}{B ^ {2}} [ B \delta A - A \delta B ] = \frac {1}{B} [ \delta A - \lambda \delta B ] =
\end{equation*}

\begin{equation*}
\begin{array}{r l} & = \frac {1}{\int \varphi^ {*} \varphi d \tau} \left[ \int \delta \varphi^ {*} H \varphi d \tau + \int \varphi^ {*} H \delta \varphi d \tau - \right. \\ & \quad \left. - \lambda \int \{\varphi^ {*} (\delta \varphi) + (\delta \varphi^ {*}) \varphi \} d \tau \right] = \\ & = \frac {1}{\int \varphi^ {*} \varphi d \tau} \left[ \int \varphi^ {*} (H - \lambda) \delta \varphi d \tau + \right. \\ & \quad \left. + \left\{\int \varphi^ {*} (H - \lambda) \delta \varphi d \tau \right\} ^ {*} \right]. \end{array}\tag{4.6}
\end{equation*}

\textbf{定理 4.1}\quad 如果函数 $\varphi$ 满足方程

\begin{equation*}
H \varphi = \lambda \varphi ,\tag{4.7}
\end{equation*}

则对任意的 $\delta \varphi$

\begin{equation*}
\delta \lambda (\varphi) = 0;\tag{4.8}
\end{equation*}

反之，如果函数满足条件(4.8)，则它必满足方程(4.7)。

\textbf{证明}\quad 如果注意到方程(4.7)，则正面结论直接由关系式(4.6)推出。现证明其逆。假定

\begin{equation*}
H \varphi - \lambda \varphi = \psi ,\tag{4.9}
\end{equation*}

这里 $\psi \neq 0$ 。进而，我们选取 $\delta \varphi = \varepsilon \psi$ 。将它代入恒等式(4.9)和(4.6)，就得

\begin{equation*}
\begin{array}{r l} \delta \lambda (\varphi) & = \frac {\int [ (\delta \varphi^ {*}) \psi + (\delta \varphi) \psi^ {*} ] d \tau}{\int \varphi^ {*} \varphi d \tau} = \\ & = \frac {\varepsilon \int [ | \psi | ^ {2} + | \psi | ^ {2} ] d \tau}{\int \varphi^ {*} \varphi d \tau} \geqslant 0. \end{array}\tag{4.10}
\end{equation*}

由(4.10)式即知，仅当 $\psi = 0$ ，即 $H\varphi = \lambda\varphi$ 时， $\delta\lambda(\varphi) = 0$ 。这就完全证明了定理。

考虑函数的增量，它有

\begin{equation*}
\delta \varphi (x) = \varepsilon \delta_ {D} (x - x _ {0})\tag{4.11}
\end{equation*}

的形式。这里用 $\delta_{D}$ 记狄拉克 $\delta$ -函数（它不致与变分 $\delta$ 相混淆）。一般说来，函数 $\varphi$ 可以是复的： $\varphi = \varphi_{1} + i\varphi_{2}$ 。我们考虑下面两种情形：

\begin{equation*}
\begin{array}{l} 1. \delta \varphi | _ {\varphi_ {2} = \text { const }} = \delta \varphi_ {1} = \varepsilon \delta_ {D} (x - x _ {0}), \\ 2. \delta \varphi | _ {\varphi_ {1} = \text { const }} = i \delta \varphi_ {2} = \varepsilon \delta_ {D} (x - x _ {0}). \end{array}
\end{equation*}

引进下面的定义：量

\begin{equation*}
\frac {\delta \lambda (\varphi)}{\delta \varphi_ {1} (x)} \equiv \frac {\delta \lambda (\varphi)}{\varepsilon}.\tag{4.12}
\end{equation*}

称为 $\lambda(\varphi)$ 关于 $\varphi_{1}$ 的泛函导数。泛函的导数通过 $\varphi$ 的变分反映出对 $\lambda(\varphi)$ 的作用，其中 $\varphi$ 的变分具有 $\varepsilon\delta_{D}(x-x_{0})$ 的形式，它可以和在空间中源所产生的场的作用相类比。

计算泛函导数：

\begin{equation*}
\begin{array}{l} \frac {\delta \lambda (\varphi)}{\delta \varphi_ {1} (x)} \Bigg | _ {\varphi_ {2} = \text {const}} = \frac {\delta \lambda (\varphi)}{\varepsilon} = \\ = \frac {1}{\varepsilon} \cdot \frac {1}{\int \varphi^ {*} \varphi d \tau} \left[ \int \varphi^ {*} (H - \lambda) \varepsilon \delta_ {D} (x - x _ {0}) d \tau + \right. \\ \left. + \int [ \varepsilon \delta_ {D} (x - x _ {0}) ] ^ {*} (H - \lambda) \varphi d \tau ] = \right. \\ = \frac {1}{\int \varphi \varphi^ {*} d \tau} [ \{(H - \lambda) \varphi \} ^ {*} + (H - \lambda) \varphi ] | _ {x = x _ {0}}. \end{array} \tag{4.13}
\end{equation*}

类似地，对于变分 $\delta\varphi_{2}=(\varepsilon/i)\delta_{D}(x-x_{0})$ 得

\begin{equation*}
\begin{array}{l} \frac {\delta \lambda (\varphi)}{\delta \varphi_ {2} (x)} \Big | _ {\varphi_ {1} = \mathrm{const}} = \\ = \frac {- i}{\int \varphi^ {*} \varphi d \tau} [ (H - \lambda) \varphi + \{(H - \lambda) \varphi \} ^ {*} ] | _ {x = x _ {0}} \end{array}\tag{4.14}
\end{equation*}

具体选择(4.11)的 $\delta\varphi$ ，使与静电学中最简单的情形相对应。如果 $\varphi$ 考虑作为电荷密度，则 $\delta\varphi = \varepsilon\delta_{D}(x - x_{0})$ 是某一瞬间在点 $x_{0}$ 处产生的，大小为 $\varepsilon$ 的电荷。如果量 $\lambda$ 是任一电荷密度分布函数（ $\lambda$ 可以是位势、场的强度等等），则 $\delta\lambda(\varphi)/\delta\varphi(x)$ 表征由于在点 $x_{0}$ 处出现电荷而引起的 $\lambda$ 的变化。

因此，由我们对泛函导数的定义推出，如果对任意的变化 $\delta \varphi(x)$ ( $x$ 属于某个区域 $R$ ), $\delta \lambda(\varphi) = 0$ ，则对 $R$ 中的所有 $x$ ，有

\begin{equation*}
\frac {\delta \lambda (\varphi)}{\delta \varphi_ {1} (x)} = \frac {\delta \lambda (\varphi)}{\delta \varphi_ {2} (x)} = 0.
\end{equation*}

\section{泛函 $\lambda(\varphi)$ 的极值性质}

设 $H$ 是厄米算子，且

\begin{equation*}
H \varphi_ {m} = \lambda_ {m} \varphi_ {m}.\tag{4.15}
\end{equation*}

将算子(4.15)特征值的次序调整为

\begin{equation*}
\lambda_ {1} \leqslant \lambda_ {2} \leqslant \lambda_ {3} \leqslant \dots ,\tag{4.16}
\end{equation*}

相应的特征函数为

\begin{equation*}
\varphi_ {1}, \varphi_ {2}, \varphi_ {3}, \dots .\tag{4.17}
\end{equation*}

\textbf{定理 4.2}\quad $\lambda(\varphi)$ 的极小值等于 $\lambda_{1}$ ，这里 $\varphi$ 为属于容许函数类中的任意函数。
\par\nopagebreak\vspace{-1.6em}
\vspace{0.15em}

\textbf{证明}\quad 假定 $\varphi_{\alpha}$ 是使泛函 $\lambda (\varphi)$ 达到相对极小值的函数，则

\begin{equation*}
\frac {\delta \lambda (\varphi)}{\delta \varphi_ {\alpha} (x)} = 0,
\end{equation*}

即
\begin{equation*}
\delta \lambda (\varphi) | _ {\varphi = {\varphi} _ {\alpha}} = 0
\end{equation*}

(根据相对极小的定义)。这表明(根据定理4.1)，当 $\varphi = \varphi_{\alpha}(x)$ 时，

\begin{equation*}
H \varphi_ {\alpha} (x) = \lambda_ {\alpha} \varphi_ {\alpha} (x).
\end{equation*}

由此即得，数 $\lambda_{\alpha}$ 是泛函 $\lambda (\varphi)$ 的极值。由此，1） $\lambda_{1}$ 是绝对极小值，因 $\lambda_{1}$ 是所有 $\lambda_{i}$ 中的最小值；2）在函数 $\varphi_{1}$ 上达到绝对极小值，定理证毕。

对于正交于 $\varphi_{1}$ 的函数 $\varphi_{i}$ 全体，则由此定理可得下面的推论。

推论 在正交于 $\varphi_{1}$ 的函数全体，即用条件

\begin{equation*}
\int \varphi_ {i} ^ {*} \varphi_ {1} d \tau = 0
\end{equation*}

所确定的函数 $\varphi_{i}$ 全体上的泛函 $\lambda(\varphi)$ ，其绝对极小值等于特征值 $\lambda_{2}$ 。

一般地，定义在正交于由 $\varphi_{1}, \varphi_{2}, \cdots, \varphi_{m-1}$ 所构成的子空间的函数全体上的泛函数 $\lambda(\varphi)$ ，其绝对极小值等于 $\lambda_{m}$ ，而达到此最小值的函数是特征函数 $\varphi_{m}$ 。

这个结论的证明可归结为如同 $\lambda_{1}$ 的情形，所不同的是，考虑的不是整个空间，而只是正交于 $\{\varphi_{1},\varphi_{2},\cdots,\varphi_{m-1}\}$ 的函数空间。

这样一来，我们就获得了逐步寻求厄米算子特征值 $\lambda_{m}$ 和特征函数 $\varphi_{m}$ 的方法，此法也适用于数值逼近方法，即先给出函数 $\varphi_{m}$ 的形式，认为它依赖于若干个参数，然后求取作为这些参数的函数的 $\lambda(\varphi_{m})$ 的极小值；或者利用某种逐次逼近的方法，找到满意的 $\varphi_{m}$ 值，因而找到满意的 $\lambda_{m}$ 值。这种方法常用在量子力学中。

\textbf{定理 4.3}\quad （完备性定理）如果最小特征值有限，且 $\lim_{n\to\infty}\lambda_{n}=+\infty$ ，则利用上面所描述的方法构造的函数系 $\{\varphi_{1},\varphi_{2},\varphi_{3},\cdots\}$ 是完备正交系。

\textbf{证明}\quad 由作法知 $\varphi_{1}, \varphi_{2}, \varphi_{3}, \cdots$ 互相正交，因此只需证明其为完备系。在作出一般情形的证明之前，先注意到我们的结论对两个特殊形式的算子成立。

容易检验，对于算子 $H = -\nabla^2$ 定理成立，此算子作用于定义域是边长为 $L$ 的立方体的函数。在这种情形， $\lambda = k_1^2 + k_2^2 + k_3^2$ ，其中每一个 $k_i$ 取值

\begin{equation*}
\frac {2 m _ {i} \pi}{L} \quad (m _ {i} = 0, 1, 2, \dots).
\end{equation*}

定理对算子 $H = L^2$ 也是成立的，此算子作用于定义在球上的函数[对于它 $\lambda = l(l + 1)]$ 。事实上在两种情形，最小特征值 $\lambda$ 都是有限的(且等于零)，此外两者的特征值 $\lambda$ 都是无界的。正如我们早已阐明的那样，这两个算子的特征函数都是完备正交系，且第一个算子是按三角函数展开，第二个算子则是按球谐函数展开。

现转到一般情形的证明。考虑函数 $f$ ，并引进下面的记号：

\begin{equation*}
\psi_ {n} \equiv f - \sum_ {i = 1} ^ {n} a _ {i} \varphi_ {i},\tag{4.18}
\end{equation*}

这里
\begin{equation*}
a _ {i} = \int \varphi_ {i} ^ {*} f d \tau .
\end{equation*}

等式(4.18)两边乘以 $\varphi_{i}^{*}$ 并积分之，得

\begin{equation*}
\int \varphi_ {i} ^ {*} \psi_ {n} d \tau = a _ {i} - a _ {i} = 0.\tag{4.19}
\end{equation*}

由此可见 $\psi_{n}$ 正交于由 $\varphi_{1}, \varphi_{2}, \cdots, \varphi_{n}$ 构成的线性子空间。即

\begin{equation*}
\psi_ {n} \perp \varphi_ {1}, \quad \varphi_ {2}, \quad \dots , \quad \varphi_ {n}.
\end{equation*}

其次考虑表示式

\begin{equation*}
\lambda (\psi_ {n}) = \frac {\int \psi_ {n} ^ {*} H \psi_ {n} d \tau}{\int \psi_ {n} ^ {*} \psi_ {n} d \tau}.\tag{4.20}
\end{equation*}

因为 $\psi_{n}\perp\varphi_{1},\varphi_{2},\cdots,\varphi_{n}$ 。故由定理 4.2 的推论，

\begin{equation*}
\min [ \lambda (\psi_ {n}) ] = \lambda_ {n + 1}.\tag{4.21}
\end{equation*}

利用(4.18)，将表达式(4.20)的分子改写成

\begin{equation*}
\begin{array}{l} \int \psi_ {n} ^ {*} H \psi_ {n} d \tau = \\ = \int \left[ f ^ {*} - \sum_ {i = 1} ^ {n} a _ {i} ^ {*} \varphi_ {i} ^ {*} \right] H \left[ f - \sum_ {i = 1} ^ {n} a _ {i} \varphi_ {i} \right] d \tau = \\ = \int f ^ {*} H f d \tau - \sum_ {i = 1} ^ {n} | a _ {i} | ^ {2} \lambda_ {i}, \end{array}\tag{4.22}
\end{equation*}

这里在逐项互乘后，用到了等式

\begin{equation*}
H \varphi_ {i} = \lambda_ {i} \varphi_ {i}, \quad a _ {i} = \int f \varphi_ {i} ^ {*} d \tau .
\end{equation*}

将(4.22)的结果代入(4.20)式：

\begin{equation*}
\lambda (\psi_ {n}) = \frac {\int f ^ {*} H f d \tau - \sum_ {i = 1} ^ {n} | a _ {i} | ^ {2} \lambda_ {i}}{\int | \psi_ {n} | ^ {2} d \tau};\tag{4.23}
\end{equation*}

注意到(4.23)，由(4.21)得

\begin{equation*}
\frac {\int f ^ {*} H f d \tau - \sum_ {i = 1} ^ {n} | a _ {i} | ^ {2} \lambda_ {i}}{\int | \psi_ {n} | ^ {2} d \tau} \geqslant \lambda_ {n + 1}\tag{4.24}
\end{equation*}

或
\begin{equation*}
\lambda_ {n + 1} \int | \psi_ {n} | ^ {2} d \tau \leqslant \int f ^ {*} H f d \tau - \sum_ {i = 1} ^ {n} | a _ {i} | ^ {2} \lambda_ {i}.\tag{4.25}
\end{equation*}

根据定理条件，存在某个 $\lambda_{m} > 0$ ，且对所有 $n > m$ ， $\lambda_{n} > 0$ 。因此略去下标 $m + 1$ 到 $n$ 的项，显然(4.25)右边部分增大了一个正数，此时应该将符号 $\leqslant$ 改为 $<$ 。这样一来

\begin{equation*}
\int | \psi_ {n} | ^ {2} d \tau <   \frac {1}{\lambda_ {n + 1}} \left[ \int f ^ {*} H f d \tau - \sum_ {i = 1} ^ {m} | a _ {i} | ^ {2} \lambda_ {i} \right],\tag{4.26}
\end{equation*}

右边方括号里的式子是不依赖于 $n$ 的有限数。其次，由于定理中 $\lim_{n\to \infty}\lambda_n = \infty$ 的条件，故当 $n$ 趋向于无限时，我们就得

\begin{equation*}
\lim _ {n \rightarrow \infty} \int | \psi_ {n} | ^ {2} d \tau = 0.\tag{4.27}
\end{equation*}

但这正好表明，函数系 $\{\varphi_{n}\}$ 在希尔伯特空间中完备。定理证毕。

\section{极值定理}

设
\begin{equation*}
\lambda (\varphi) = \frac {\int_ {R} \varphi^ {*} H \varphi d \tau}{\int_ {R} \varphi^ {*} \varphi d \tau},\tag{4.28}
\end{equation*}

但现在函数 $\varphi$ 定义在区域 $R$ 内，$R$ 不一定是正规的几何图形。取任意的正交函数集合 $f_{1}, f_{2}, \cdots, f_{n}$ ，它们定义在 $R$ 上且属于容许函数类。引入记号

\begin{equation*}
m \left(f _ {1}, f _ {2}, \dots , f _ {n}\right) = \min [ \lambda (\varphi) ]\tag{4.29}
\end{equation*}

这里所有的 $\varphi$ 为正交于张在 $f_{1}, f_{2}, \cdots, f_{n}$ 上的子空间中的函数，即满足等式

\begin{equation*}
\int \varphi^ {*} f _ {i} d \tau = 0
\end{equation*}

的函数 $\varphi$ 。于是有下面的定理。

\textbf{定理 4.4}\quad 下面的等式成立：

\begin{equation*}
\max \left\{m (f _ {1}, f _ {2}, \dots , f _ {n}) \right\} = \lambda_ {n + 1}.\tag{4.30}
\end{equation*}

\textbf{证明}\quad 象通常那样， $H\varphi_{i} = \lambda_{i}\varphi_{i}$ 。记

\begin{equation*}
\varphi = \sum_ {i = 1} ^ {n + 1} a _ {i} \varphi_ {i},\tag{4.31}
\end{equation*}

并要求 $\varphi$ 正交于由 $f_{1}, f_{2}, \cdots, f_{n}$ 构成的子空间。这个要求通常能满足，因为总可选取 $n+1$ 个系数，使之满足 $n$ 个条件。因为 [根据 (4.31)]

\begin{equation*}
H \varphi = \sum_ {i = 1} ^ {n + 1} a _ {i} \lambda_ {i} \varphi_ {i},
\end{equation*}

故
\begin{equation*}
\int \varphi^ {*} H \varphi d \tau = \sum_ {i = 1} ^ {n + 1} | a _ {i} | ^ {2} \lambda_ {i}\tag{4.32}
\end{equation*}

和
\begin{equation*}
\int \varphi^ {*} \varphi d \tau = \sum_ {i = 1} ^ {n + 1} | a _ {i} | ^ {2}.\tag{4.33}
\end{equation*}

将等式(4.32)和(4.33)代入(4.28)式，得

\begin{equation*}
\lambda (\varphi) = \frac {\sum_ {i = 1} ^ {n + 1} \left| a _ {i} \right| ^ {2} \lambda_ {i}}{\sum_ {i = 1} ^ {n + 1} \left| a _ {i} \right| ^ {2}}.\tag{4.34}
\end{equation*}

所以 $\lambda (\varphi)$ 是 $\lambda_{i}$ 的加权平均，它不可能超过 $\lambda_{i}$ 中的最大者。因为所有 $\lambda_{i} \leqslant \lambda_{n + 1}$ （当 $i \leqslant n + 1$ ），故

\begin{equation*}
\lambda (\varphi) = \frac {\sum_ {i = 1} ^ {n + 1} | a _ {i} | ^ {2} \lambda_ {i}}{\sum_ {i = 1} ^ {n + 1} | a _ {i} | ^ {2}} \leqslant \lambda_ {n + 1},\tag{4.35}
\end{equation*}

因此若注意到(4.29)，则有

\begin{equation*}
m (f _ {1}, f _ {2}, \dots , f _ {n}) \leqslant \lambda (\varphi) \leqslant \lambda_ {n + 1}.\tag{4.36}
\end{equation*}

于是我们看到 $m(f_{1}, f_{2}, \cdots, f_{n})$ 的最大值不能大于 $\lambda_{n+1}$ 。现在证明， $m(f_{1}, f_{2}, \cdots, f_{n})$ 的最大值达到值 $\lambda_{n+1}$ 。为此，我们研究 $(f_{1}, f_{2}, \cdots, f_{n}) = (\varphi_{1}, \varphi_{2}, \cdots, \varphi_{n})$ 的特殊情形。由排列特征值的方法本身推出

\begin{equation*}
m \left(\varphi_ {1}, \varphi_ {2}, \dots , \varphi_ {n}\right) = \lambda_ {n + 1},\tag{4.37}
\end{equation*}

因此在一般情形
